\documentclass[11pt]{article}
\usepackage[a4paper,margin=1in]{geometry}
\usepackage{booktabs}
\usepackage{longtable}
\usepackage{array}
\usepackage{hyperref}
\usepackage{graphicx}
\usepackage{caption}
\usepackage{amsmath}
\usepackage[T1]{fontenc}
\usepackage{lmodern}

% Explicit run folder used for embedded evidence figures.
\edef\RunDir{\detokenize{../run_outputs/20260719_191707}}

\title{ToyModelSOEMC Training Run: Explanation of All MATLAB Outputs}
\author{Training\_2x2x2\_share package}
\date{}

\begin{document}
\maketitle

\section{Purpose}
This document explains every output produced by the training run script
\texttt{ToyModelSOEMC\_run\_training\_2x2x2.m}, which launches
\texttt{ToyModelSOEMC\_run.m} with:
\begin{itemize}
\item \texttt{MINIMAL\_MODE = false}
\item \texttt{TRAINING\_MODE = true}
\end{itemize}

In this configuration, the joint policy search in Step 4 uses a reduced
$2 \times 2 \times 2$ grid:
\begin{itemize}
\item $d \in \{90, 365\}$ days
\item $\phi_R \in \{0, 1.0\}$
\item $\phi_{IQ} \in \{0.04, 0.08\}$
\end{itemize}

\section{Explicit Run Outputs Used in This Report}
This report is based on the concrete exported artifacts from the run folder:

\begin{itemize}
\item \texttt{run\_outputs/20260719\_191707/}
\item Manifest timestamp: \texttt{2026-07-19 19:20:14}
\item Run success flag in manifest: \texttt{1}
\end{itemize}

The six exported figure sets detected in that folder are:
\begin{enumerate}
\item \texttt{01\_Supply\_Crisis\_Demo.(png|pdf|fig)}
\item \texttt{02\_Vietnam\_Reserve\_Adequacy\_Under\_Growth.(png|pdf|fig)}
\item \texttt{03\_Energy\_Master\_Plan\_VIII\_Coal\_Transition.(png|pdf|fig)}
\item \texttt{04\_Joint\_Optimal\_Policy\_Heatmap\_at\_d.(png|pdf|fig)}
\item \texttt{05\_Joint\_Optimal\_Welfare\_vs\_Reserve\_Level.(png|pdf|fig)}
\item \texttt{06\_Joint\_Optimal\_Stock-Drawdown\_Interaction.(png|pdf|fig)}
\end{enumerate}

In addition, the policy summary values below are taken from
\texttt{joint\_optimal\_results.txt} generated by the same package run.

\section{Observed Numerical Outputs (from joint\_optimal\_results.txt)}
\begin{center}
\begin{tabular}{ll}
\toprule
Metric & Value \\
\midrule
$d^*$ (days) & 90 \\
$\phi_R^*$ & 0.50 \\
$\phi_{IQ}^*$ & 0.04 \\
$R_{oil}^*$ & 0.0740 \\
$R_{coal}^*$ & 0.1110 \\
$E[W]_{opt}$ & -44.1433 \\
$E[W]_{zero}$ & -44.1433 \\
Storage cost (\% of $C$) & 4.444 \\
$\lambda_{benefit}$ (\%) & 0.000 \\
$C_{ss}$ & 0.1313 \\
$Y_{ss}$ & 0.1626 \\
$Q_{oil}^*$ & 0.0650 \\
$Q_{coal}^*$ & 0.0650 \\
$N_{MC}$ & 40 \\
$N_d$ & 2 \\
\bottomrule
\end{tabular}
\end{center}

\section{Output Inventory (At a Glance)}
The MATLAB run produces four categories of outputs:
\begin{enumerate}
\item Console output (progress, diagnostics, policy summaries)
\item Interactive figures (scenario and policy plots)
\item Saved data/text artifacts (result files)
\item Workspace state updates (updated model parameters and arrays)
\end{enumerate}

\section{Console Outputs and Meaning}
\subsection*{Step 1: Dynare compilation and initialization}
\textbf{Printed header:} \texttt{=== STEP 1: Running Dynare ===}

\textbf{What it means:}
\begin{itemize}
\item Dynare parses \texttt{ToyModelSOEMC.mod}.
\item Model objects \texttt{M\_}, \texttt{options\_}, \texttt{oo\_} are created.
\item The custom steady-state file \texttt{ToyModelSOEMC\_steadystate.m} is used.
\end{itemize}

\subsection*{Step 2: Steady-state check}
\textbf{Printed header:} \texttt{=== STEP 2: Steady-State Check (Vietnam calibration) ===}

\textbf{Key printed values:}
\begin{itemize}
\item $Y_{ss}$, $C_{ss}$, $r_{ss}$, $E_{ss}$, $NX_{ss}$
\item Ratios such as $C_{ss}/Y_{ss}$ and $I_{ss}/Y_{ss}$
\item Commodity-level reserve and capacity diagnostics
\end{itemize}

\textbf{How to interpret:}
\begin{itemize}
\item Confirms calibration consistency before policy simulations.
\item Checks whether steady-state moments are plausible for the intended Vietnam setting.
\end{itemize}

\subsection*{Step 3: Deterministic crisis demo}
\textbf{Printed header:} \texttt{=== STEP 3: Deterministic Crisis Demo ===}

\textbf{Key printed lines:}
\begin{itemize}
\item Crisis specification (supply disruption path)
\item Scenario welfare values $W$ for each $\phi_R$ scenario
\end{itemize}

\textbf{How to interpret:}
\begin{itemize}
\item Demonstrates reserve drawdown policy effect under a known shock path.
\item The relative $W$ values summarize which drawdown rule performs better in that deterministic setup.
\end{itemize}

\subsection*{Step 3b: Growth adequacy scenario}
\textbf{Printed header:} \texttt{=== STEP 3b: Vietnam Growth Scenario ===}

\textbf{Key printed lines:}
\begin{itemize}
\item Current reserve coverage (days)
\item IEA benchmark (90 days)
\item End-horizon coverage under fixed reserves
\end{itemize}

\textbf{How to interpret:}
\begin{itemize}
\item Shows that with GDP/import growth, fixed nominal reserves imply declining coverage ratios.
\item Supports policy logic of scaling reserve targets with demand.
\end{itemize}

\subsection*{Step 3c: Energy Master Plan transition scenarios}
\textbf{Printed header:} \texttt{=== STEP 3c: Energy Master Plan VIII --- Coal Transition Scenario ===}

\textbf{Key printed lines:}
\begin{itemize}
\item Scenario labels and welfare values $W$
\item Qualitative interpretation block (A/B/C scenario comparison)
\end{itemize}

\textbf{How to interpret:}
\begin{itemize}
\item Compares no-policy, investment, and reserve-buffer responses during coal phase-down.
\item Highlights interaction between investment and reserves as transition tools.
\end{itemize}

\subsection*{Step 4: Joint optimization}
\textbf{Printed header:} \texttt{=== STEP 4: Joint Optimal Policy (phi\_R x phi\_IQ x d) ===}

\textbf{Key printed lines:}
\begin{itemize}
\item Active policy grids and Monte Carlo count
\item Per-grid-point expected welfare values
\item Final optimal tuple $(d^*, \phi_R^*, \phi_{IQ}^*)$
\end{itemize}

\textbf{How to interpret:}
\begin{itemize}
\item This is the central optimization result for policy recommendation.
\item In training mode, result quality is illustrative because grid and run count are reduced.
\end{itemize}

\section{Figures Produced by MATLAB and Interpretation}
\begin{longtable}{>{\raggedright\arraybackslash}p{0.28\textwidth} >{\raggedright\arraybackslash}p{0.68\textwidth}}
\toprule
\textbf{Figure name (window title)} & \textbf{What it shows and how to read it} \\
\midrule
\texttt{Supply Crisis Demo} & Time paths for consumption, energy services, and reserve stocks under deterministic disruption. Compare scenario lines to see buffering strength of policy settings. \\
\texttt{Vietnam: Reserve Adequacy Under Growth} & Coverage (days) over time as demand grows. The gap between fixed-reserve line and target lines visualizes adequacy erosion if stock does not scale. \\
\texttt{Energy Master Plan VIII --- Coal Transition} & Consumption, energy-service, and coal-stock trajectories for transition scenarios A/B/C. Quantifies trade-offs among no response, higher investment, and reserve use. \\
\texttt{Joint Optimal: Policy Heatmap at d*} & Heatmap of expected welfare at optimal reserve level across $(\phi_R, \phi_{IQ})$. Brightest region marks best policy combination for that $d$. \\
\texttt{Joint Optimal: Welfare vs Reserve Level} & Expected welfare as a function of reserve days at optimal $(\phi_R^*, \phi_{IQ}^*)$. The peak indicates preferred reserve sizing. \\
\texttt{Joint Optimal: Stock-Drawdown Interaction} & Welfare-versus-$d$ curves for different $\phi_R$ values, showing complementarity between stock size and drawdown aggressiveness. \\
\bottomrule
\end{longtable}

\section{Embedded Output Figures (Explicitly Included)}
This section embeds the concrete figure exports from
the run output folder \texttt{20260719\_191707} directly into the presentation document.

\subsection*{Figure 01: Supply Crisis Demo}
\begin{figure}[h!]
\centering
\IfFileExists{\RunDir/01_Supply_Crisis_Demo.png}{
	\includegraphics[width=0.95\textwidth]{\RunDir/01_Supply_Crisis_Demo.png}
}{
	\fbox{\parbox{0.9\textwidth}{Missing file: \texttt{01\_Supply\_Crisis\_Demo.png}}}
}
\caption{Deterministic crisis comparison of policy responses.}
\end{figure}

\subsection*{Figure 02: Reserve Adequacy Under Growth}
\begin{figure}[h!]
\centering
\IfFileExists{\RunDir/02_Vietnam_Reserve_Adequacy_Under_Growth.png}{
	\includegraphics[width=0.95\textwidth]{\RunDir/02_Vietnam_Reserve_Adequacy_Under_Growth.png}
}{
	\fbox{\parbox{0.9\textwidth}{Missing file: \texttt{02\_Vietnam\_Reserve\_Adequacy\_Under\_Growth.png}}}
}
\caption{Coverage erosion under GDP/import growth with fixed reserves.}
\end{figure}

\subsection*{Figure 03: Energy Master Plan VIII Transition}
\begin{figure}[h!]
\centering
\IfFileExists{\RunDir/03_Energy_Master_Plan_VIII_Coal_Transition.png}{
	\includegraphics[width=0.95\textwidth]{\RunDir/03_Energy_Master_Plan_VIII_Coal_Transition.png}
}{
	\fbox{\parbox{0.9\textwidth}{Missing file: \texttt{03\_Energy\_Master\_Plan\_VIII\_Coal\_Transition.png}}}
}
\caption{Coal phase-down scenarios and transition dynamics.}
\end{figure}

\subsection*{Figure 04: Joint Optimal Heatmap}
\begin{figure}[h!]
\centering
\IfFileExists{\RunDir/04_Joint_Optimal_Policy_Heatmap_at_d.png}{
	\includegraphics[width=0.95\textwidth]{\RunDir/04_Joint_Optimal_Policy_Heatmap_at_d.png}
}{
	\fbox{\parbox{0.9\textwidth}{Missing file: \texttt{04\_Joint\_Optimal\_Policy\_Heatmap\_at\_d.png}}}
}
\caption{Policy interaction heatmap at the selected reserve level.}
\end{figure}

\subsection*{Figure 05: Welfare vs Reserve Level}
\begin{figure}[h!]
\centering
\IfFileExists{\RunDir/05_Joint_Optimal_Welfare_vs_Reserve_Level.png}{
	\includegraphics[width=0.95\textwidth]{\RunDir/05_Joint_Optimal_Welfare_vs_Reserve_Level.png}
}{
	\fbox{\parbox{0.9\textwidth}{Missing file: \texttt{05\_Joint\_Optimal\_Welfare\_vs\_Reserve\_Level.png}}}
}
\caption{Expected welfare profile over reserve coverage days.}
\end{figure}

\subsection*{Figure 06: Stock--Drawdown Interaction}
\begin{figure}[h!]
\centering
\IfFileExists{\RunDir/06_Joint_Optimal_Stock-Drawdown_Interaction.png}{
	\includegraphics[width=0.95\textwidth]{\RunDir/06_Joint_Optimal_Stock-Drawdown_Interaction.png}
}{
	\fbox{\parbox{0.9\textwidth}{Missing file: \texttt{06\_Joint\_Optimal\_Stock-Drawdown\_Interaction.png}}}
}
\caption{Complementarity between reserve stock size and drawdown aggressiveness.}
\end{figure}

\clearpage

\section{Saved Files and Their Content}
\subsection*{\texttt{joint\_optimal\_results.txt}}
Produced by \texttt{ToyModelSOEMC\_joint\_optimal.m}. Contains:
\begin{itemize}
\item Optimal policy tuple: \texttt{d\_star}, \texttt{phi\_R\_star}, \texttt{phi\_IQ\_star}
\item Implied reserve stocks and days coverage
\item Welfare metrics: \texttt{EW\_opt}, \texttt{EW\_zero}
\item Storage-cost and welfare-benefit summary fields
\item Calibration and run metadata (e.g., \texttt{N\_MC}, \texttt{N\_d})
\end{itemize}

For this explicit run, the file reports:
\begin{itemize}
\item \texttt{d\_star=90}, \texttt{phi\_R\_star=0.50}, \texttt{phi\_IQ\_star=0.04}
\item \texttt{N\_MC=40}, \texttt{N\_d=2}
\end{itemize}

\subsection*{\texttt{ToyModelSOEMC\_steadystate.mat}}
Produced by \texttt{ToyModelSOEMC\_steadystate.m}. Stores:
\begin{itemize}
\item \texttt{ys}: steady-state vector in model variable order
\item \texttt{params}: updated parameter vector (including capacity-bar updates)
\end{itemize}

\subsection*{Complete output-file glossary (clear explanations)}
\small
\begin{longtable}{>{\raggedright\arraybackslash}p{0.38\textwidth} >{\raggedright\arraybackslash}p{0.57\textwidth}}
\toprule
\textbf{File} & \textbf{Purpose} \\
\midrule
\path{joint_optimal_results.txt} & Plain-text summary of optimal policy and welfare metrics; primary numeric output for reporting. \\
\path{ToyModelSOEMC_steadystate.mat} & MATLAB binary with solved steady state and updated parameter vector. \\
\path{ToyModelSOEMC.log} & Dynare execution log for debugging preprocessing/solve issues. \\
\path{run_outputs/<timestamp>/matlab_console.log} & Captured console output from export run (diary), if available. \\
\path{run_outputs/<timestamp>/export_manifest.txt} & Index of all exported figure files and their paths, with timestamp and success flag. \\
\path{run_outputs/<timestamp>/NN_<name>.png} & Raster figure export for slides/documents. \\
\path{run_outputs/<timestamp>/NN_<name>.pdf} & Vector figure export for high-quality print/LaTeX use. \\
\path{run_outputs/<timestamp>/NN_<name>.fig} & Native editable MATLAB figure format. \\
\path{run_outputs/<timestamp>/workspace_snapshot.mat} & End-of-run workspace snapshot for reproducibility and post-run analysis. \\
\path{ToyModelSOEMC_outputs_explained.tex} & Source of this documentation report. \\
\path{ToyModelSOEMC_outputs_explained.pdf} & Compiled documentation with embedded output figures. \\
\path{ToyModelSOEMC_outputs_explained.aux/.log/.out/.fls/.fdb_latexmk/.synctex.gz} & LaTeX intermediate/build files; safe to regenerate. \\
\bottomrule
\end{longtable}
\normalsize

\section{Workspace Outputs (Not Auto-Saved)}
After run completion, MATLAB workspace includes key variables:
\begin{itemize}
\item From run script: scenario paths/structures (e.g., \texttt{results\_demo}, \texttt{results\_emp})
\item From joint optimization: \texttt{W\_joint}, \texttt{phi\_R\_star}, \texttt{phi\_IQ\_star}, \texttt{d\_star}
\item Dynare objects: \texttt{M\_}, \texttt{oo\_}, \texttt{options\_}
\end{itemize}

These are useful for additional analysis but are volatile unless explicitly saved.

\section{Why Results Can Differ Across Runs}
Some outputs are stochastic because Step 4 uses Monte Carlo shocks in PF mode.
Differences arise from:
\begin{itemize}
\item Random shock draws in \texttt{ToyModelSOEMC\_shock\_gen.m}
\item Solver convergence path (mitigated with homotopy steps)
\item Selected run mode (full vs training vs minimal)
\end{itemize}

\section{Recommended Reporting Practice}
For reproducible reporting, save these artifacts after each run:
\begin{enumerate}
\item Console log (MATLAB \texttt{diary})
\item Figure exports (PNG/PDF)
\item \texttt{joint\_optimal\_results.txt}
\item \texttt{ToyModelSOEMC\_steadystate.mat}
\end{enumerate}

This document can then be paired with the exported figures for a complete training report.

\section{Note on Console Log Availability}
In the explicit run folder used above, \texttt{matlab\_console.log} exists but is empty.
The figure and policy output files are still valid and were used as primary evidence.

\end{document}
